\chapter{赋范向量空间}\label{ch:b2:nvs}

当度量空间是向量空间，且距离来自
范数时，拓扑与线性代数开始相互作用：线性映射
恰在单位球上有界时连续，有限维空间迫使所有范数一致，
完备性使绝对收敛的级数成为收敛级数。
有限维与无限维的分界，
在里斯定理中得以明确，
是本章最深刻的教训。

下文中，设 $E, F$ 为 $K = \R$ 或 $\C$ 上的向量空间。

\section{范数}

\begin{definition}\label{def:b2:nvs:norm}
称\emph{范数}\index{范数}为 $E$ 上的映射 $\norm{\,\cdot\,} \colon E
\to \R_+$，且对所有 $x, y \in E$、$\lambda \in K$ 满足：
\[
\norm x = 0 \iff x = 0,
\qquad
\norm{\lambda x} = \abs\lambda\,\norm x,
\qquad
\norm{x + y} \leq \norm x + \norm y .
\]
于是 $d(x, y) = \norm{x - y}$ 为距离，
\cref{ch:b2:metric} 中的全部结论均适用。
逆三角不等式 $\bigl|\norm x - \norm y\bigr| \leq \norm{x - y}$ 表明
范数本身是 $1$-利普希茨的；加法与数乘
是连续的（估计式 $\norm{(x + y) - (x' + y')} \leq \norm{x -
x'} + \norm{y - y'}$ 等）。
\end{definition}

\begin{example}\label{ex:b2:nvs:examples}
在 $K^n$ 上：
\[
\norm{x}_1 = \sum_i \abs{x_i},
\qquad
\norm{x}_2 = \Bigl(\sum_i \abs{x_i}^2\Bigr)^{1/2},
\qquad
\norm{x}_\infty = \max_i \abs{x_i}
\]
（$\norm\cdot_2$ 是范数，由柯西–施瓦茨不等式得出，见本科一年级分册）。
在 $C(\intcc{a}{b})$ 上：
\[
\norm f_\infty = \sup \abs f,
\qquad
\norm f_1 = \int_a^b \abs f,
\qquad
\norm f_2 = \Bigl(\int_a^b \abs f^2\Bigr)^{1/2},
\]
后两者之所以是范数，依赖于积分的严格正性，
以及积分形式的柯西–施瓦茨不等式（见本科一年级分册）。对矩阵而言：
$\mathcal M_n(K) \simeq K^{n^2}$ 上的任一范数；
下面的算子范数在结构上最为重要。
\end{example}

\begin{definition}[等价范数]\label{def:b2:nvs:equivalent}
设 $N_1, N_2$ 为 $E$ 上的两个范数，称它们为\emph{等价范数}\index{等价范数}，
若存在常数 $c, C > 0$，使得
\[
c\,N_1 \leq N_2 \leq C\, N_1 .
\]
等价范数具有相同的开集、相同的收敛序列与柯西序列，
相同的紧子集与完备子集：
即同一套分析。
\end{definition}

\begin{example}[无限维中的不等价]\label{ex:b2:nvs:nonequivalent}
在 $C(\intcc{0}{1})$ 上，总有 $\norm f_1 \leq \norm f_\infty$，
但反向估计不成立：$f_n(x) = x^n$ 满足 $\norm{f_n}_\infty = 1$，
而 $\norm{f_n}_1 = \frac{1}{n+1} \to 0$。因此 $f_n \to 0$
在 $\norm\cdot_1$ 下成立，而在 $\norm\cdot_\infty$ 下不成立：两个范数
对“收敛”本身的判断并不一致。
\end{example}

\begin{example}[$n$ 维中的显式常数]\label{ex:b2:nvs:constants}
在 $K^n$ 上，三个经典范数等价，且
可取最优常数：
\[
\norm x_\infty \leq \norm x_2 \leq \norm x_1
\leq \sqrt n\,\norm x_2 \leq n\,\norm x_\infty ,
\]
中间的不等式 $\norm x_1 \leq \sqrt n\norm x_2$ 由
对全 1 向量使用柯西–施瓦茨不等式得到。取等号的向量：
$e_1$ 使前两个不等式取等号，$(1, 1,
\dots, 1)$ 使后两个取等号。维数 $n$ 清楚地体现在
常数之中，这便是无限维情形失效的定量根源：
当 $n \to \infty$ 时，任何一致的常数都不复存在，
这正是 \cref{ex:b2:nvs:nonequivalent} 所展示的现象，
即发生在函数空间上的情形。
\end{example}

\begin{omfigure}
\begin{tikzpicture}[scale=1.5]
  \draw[->] (-1.45,0) -- (1.5,0) node[below, font=\small] {$x$};
  \draw[->] (0,-1.45) -- (0,1.5) node[left, font=\small] {$y$};
  \draw[omProp, very thick]
    (1,1) -- (-1,1) -- (-1,-1) -- (1,-1) -- cycle;
  \draw[omThm, very thick] (0,0) circle (1);
  \draw[omDef, very thick]
    (1,0) -- (0,1) -- (-1,0) -- (0,-1) -- cycle;
  \node[omDef, font=\small] at (0.36,0.28)
    {$\norm\cdot_1$};
  \node[omThm, font=\small] at (0.55,0.92)
    {$\norm\cdot_2$};
  \node[omProp, font=\small] at (1.24,1.13)
    {$\norm\cdot_\infty$};
\end{tikzpicture}

{\small $\R^2$ 上三个经典范数的单位球，
按 \cref{ex:b2:nvs:constants} 中的不等式嵌套排列，
即球越小、范数越大。球的圆润性至关重要：菱形与正方形的平直边
恰是严格凸性不成立之处，
这一性质在 \cref{ch:b2:realfun} 中的周末问题
及本章第 4 题中被利用。}
\end{omfigure}

\section{连续线性映射}

\begin{theorem}[刻画]\label{thm:b2:nvs:continuouslinear}
设 $u \colon E \to F$ 为赋范空间之间的线性映射，
以下各条件等价：
\begin{enumerate}
  \item $u$ 连续；
  \item $u$ 在 $0$ 处连续；
  \item $u$ 在闭单位球上有界：$\sup_{\norm x \leq
        1} \norm{u(x)} < \infty$；
  \item 存在 $C \geq 0$，使得 $\norm{u(x)} \leq C \norm x$ 对
        所有 $x$ 成立；
  \item $u$ 是利普希茨映射。
\end{enumerate}
使上述条件成立的最小 $C$ 称为\emph{算子范数}\index{算子范数}，
$\vertiii{u} = \sup_{\norm x \leq 1}\norm{u(x)} = \sup_{x \neq 0}
\frac{\norm{u(x)}}{\norm x}$；它使连续线性映射的空间 $\mathcal{L}_c(E,
F)$ 成为赋范空间，满足
\[
\vertiii{v \circ u} \leq \vertiii v\, \vertiii u .
\]
\end{theorem}

\begin{proof}
(1 $\Rightarrow$ 2) 显然。(2 $\Rightarrow$ 3)：由 $0$ 处的连续性，
取 $\varepsilon = 1$，得 $\delta$，使得 $\norm x \leq \delta
\Rightarrow \norm{u(x)} \leq 1$；由齐次性，对任一 $x$，
只要 $\norm x \leq 1$，便可先将其缩入该球再还原：
\[
\norm{u(x)} = \frac1\delta\,\norm{u(\delta x)} \leq
\frac1\delta ,
\]
由于 $\norm{\delta x} \leq \delta$。(3 $\Rightarrow$ 4)：
对 $x \neq 0$，将该界用于 $\frac{x}{\norm x}$。(4
$\Rightarrow$ 5)：$\norm{u(x) - u(y)} = \norm{u(x - y)} \leq
C\norm{x - y}$。(5 $\Rightarrow$ 1) 已知成立。

$\vertiii\cdot$ 的范数公理中，齐次性与分离性
显然成立（由 $\vertiii u = 0$ 知 $u = 0$ 在球上，从而
处处为零）；三角不等式由 $\norm{(u + v)(x)} \leq
\norm{u(x)} + \norm{v(x)}$ 得到。次乘性：$\norm{v(u(x))}
\leq \vertiii v\,\norm{u(x)} \leq \vertiii v \vertiii u \norm x$。
\end{proof}

\begin{example}\label{ex:b2:nvs:operatornorms}
在 $\bigl(C(\intcc{0}{1}), \norm\cdot_\infty\bigr)$ 上：求值映射
$f \mapsto f(0)$ 的算子范数为 $1$；积分映射 $f \mapsto
\int_0^1 f$ 的范数为 $1$；映射 $f \mapsto \int_0^1 t f(t)\dd t$
的范数为 $\int_0^1 t\,\dd t = \frac12$（上界由积分的三角不等式给出，
在 $f \equiv 1$ 时取到）。但
求导映射，从 $(C^1, \norm\cdot_\infty)$ 到 $(C^0,
\norm\cdot_\infty)$ 是\emph{不}连续的：$\norm{\sin(nx)}_\infty
= 1$，而导数的上确界范数为 $n$。线性
并不蕴含无穷维中的连续性。
\end{example}

\begin{example}[两个范数，对同一数列的两种判定]\label{ex:b2:nvs:sqrtnxn}
在 $C(\intcc01)$ 上，设 $g_n(x) = \sqrt{n}\,x^n$。则
\[
\norm{g_n}_1 = \frac{\sqrt n}{n + 1} \longrightarrow 0,
\qquad
\norm{g_n}_2^2 = \frac{n}{2n + 1} \longrightarrow \frac12,
\qquad
\norm{g_n}_\infty = \sqrt n \longrightarrow \infty :
\]
一个数列，三种范数，三种表现 --- 收敛到
零、不收敛（范数稳定在 $\frac1{\sqrt2}$，而
逐点极限为 $0$），以及趋于无穷。质量集中在
$x = 1$ 附近对 $\norm\cdot_1$ 不可见，
对 $\norm\cdot_2$ 半可见，对 $\norm\cdot_\infty$ 则占主导。在无穷维中，
“它是否收敛？”并不是关于某个数列的问题，
而是关于一个数列\emph{与}一个范数的问题。
\end{example}

\begin{method}[计算算子范数]\label{met:b2:nvs:opnorm}
总分两步进行。\emph{上界：}估计
$\norm{u(x)}$ 以 $C\norm x$ 为界，利用三角不等式、
柯西–施瓦茨不等式，或积分估计 --- 由此证明 $\vertiii u
\leq C$。\emph{见证：}给出某个 $x_0 \neq 0$
使 $\norm{u(x_0)} = C\norm{x_0}$（界可取到），或
给出一列单位向量 $x_n$，使 $\norm{u(x_n)} \to C$
（界被逼近）。两步缺一不可：仅有上界只得到 $\vertiii u \leq C$，
仅有见证只得到 $\vertiii u \geq C$。在无穷维中，
见证可能须是一个数列 --- 上确界
未必能取到
（见\cref{exo:b2:nvs:8}）。
\end{method}

\begin{example}[对角算子对各范数一视同仁]\label{ex:b2:nvs:diagonalnorm}
设 $D = \operatorname{diag}(d_1, \dots, d_n)$ 为 $K^n$ 上的对角算子，
赋以\emph{任一}范数 $\norm\cdot_1, \norm\cdot_2,
\norm\cdot_\infty$：由 $\abs{d_ix_i} \leq
\bigl(\max_j\abs{d_j}\bigr)\abs{x_i}$ 逐分量地，
$\norm{Dx} \leq \max_j\abs{d_j}\,\norm x$；而 $x = e_{j_0}$
（一个取最大值的指标）取到该界。于是 $\vertiii D =
\max_j\abs{d_j}$，三种情形皆然。对对角映射而言，所有
合理范数讲的是同一回事：最大的伸缩
因子。算子范数的一切难点，都在于
\emph{非}对角行为 --- 这正是本章周末问题（第 22 题）中的
适应范数借助先迫使矩阵化为对角形式而奏效的
原因。
\end{example}

\begin{example}[列和：\cref{exo:b2:nvs:4} 的 $1$-范数版本]\label{ex:b2:nvs:columnsums}
在 $(\R^n, \norm\cdot_1)$ 上，矩阵 $A$ 的算子范数是
绝对值的\emph{列}和中的最大者。依方法计算：对
$\norm x_1 \leq 1$，
\[
\norm{Ax}_1 = \sum_i\Bigl|\sum_j a_{ij}x_j\Bigr|
\leq \sum_j \abs{x_j}\sum_i\abs{a_{ij}}
\leq \Bigl(\max_j\sum_i\abs{a_{ij}}\Bigr)\norm x_1 ,
\]
且界在 $x = e_{j_0}$ 处取到，其中 $j_0$ 为取最大值的
列 --- 可想到的最整洁的见证。于是对 $A =
\begin{psmallmatrix}1 & -2\\ 3 & 1\end{psmallmatrix}$：
$\vertiii A_1 = \max(1 + 3,\ 2 + 1) = 4$，而
$\vertiii A_\infty = 4$ 同样成立（按行）--- 这里只是巧合，并非
一条规律：若不对称地调换矩阵元素，两种范数
便分道扬镳。对 $\norm\cdot_\infty$ 按行，对
$\norm\cdot_1$ 按列。助记的办法是：各范数的单位向量
（分别为符号模式与基向量）挑出相应的
和。
\end{example}

\begin{proposition}[双线性映射]\label{prop:b2:nvs:bilinear}
双线性映射 $b \colon E \times F \to G$ 连续，当且仅当
$\norm{b(x,y)} \leq C\norm x\,\norm y$ 对某个 $C$ 成立；此时它
在有界集上是利普希茨的。（证明方法相同；乘积映射
$(u, v) \mapsto v \circ u$ 与矩阵乘法是其中的关键
例子。）
\end{proposition}

\begin{proof}
若上述估计成立：
\[
b(x,y) - b(x_0,y_0) = b(x - x_0,\, y) + b(x_0,\, y - y_0),
\]
于是 $\norm{b(x,y) - b(x_0,y_0)} \leq C\norm{x - x_0}\norm y +
C\norm{x_0}\norm{y - y_0}$，即在 $(x_0, y_0)$ 处连续，并有
利普希茨型估计，其中 $\norm x, \norm y \leq R$。反之，
由 $(0,0)$ 处的连续性，存在 $\delta$，使 $\norm{b(x,y)} \leq 1$ 在
$\norm x, \norm y \leq \delta$ 上成立；对两个变量同时作伸缩。
\end{proof}

\section{有限维}

\begin{theorem}[有限维空间中范数的等价性]\label{thm:b2:nvs:finitedim}
在有限维空间上，\emph{所有范数彼此
等价}。\index{范数的等价性}因此在有限维中：
收敛性、开性、紧性、完备性都是
与范数无关的概念；紧 $=$ 闭且有界；
空间是完备的；并且每个线性（或多重线性）映射
\emph{自}有限维空间出发时必连续。
\end{theorem}

\begin{proof}
取定一组基，视 $E \simeq K^n$；只需将任一
范数 $N$ 与 $\norm\cdot_\infty$ 比较即可。

\emph{一个方向是代数的：}$N(x) = N(\sum x_i e_i) \leq
\sum \abs{x_i} N(e_i) \leq C \norm x_\infty$，其中 $C = \sum N(e_i)$。
由此还可知 $N$ 在 $(K^n, \norm\cdot_\infty)$ 上连续
（事实上它是 $C$-利普希茨的：$\abs{N(x) - N(y)} \leq N(x - y)$）。

\emph{另一个方向是拓扑的：}单位球面 $S = \{x :
\norm{x}_\infty = 1\}$ 在 $(K^n,
\norm\cdot_\infty)$ 中是闭且有界的，因而紧
（见\cref{thm:b2:metric:compactprops}~(2)，对 $\C^n \simeq
\R^{2n}$ 亦成立）。连续函数 $N$ 的最小值 $c$ 在
$S$ 上取得；$c > 0$，因为 $N$ 仅在 $0 \notin S$ 处为零。
由齐次性，该界推广为：$N(x) \geq c \norm{x}_\infty$ 对一切 $x$ 成立。

由此，所有结论都化归为 $(K^n,
\norm\cdot_\infty)$ 上的情形，而此时它们已知
（\cref{thm:b2:metric:rncomplete}，
\cref{thm:b2:metric:compactprops}）；设线性映射 $u$ 自
有限维 $E$ 出发，满足 $\norm{u(x)} \leq \sum\abs{x_i}
\norm{u(e_i)} \leq C'\norm{x}_\infty$，此即
\cref{thm:b2:nvs:continuouslinear}中的估计 (4)。
\end{proof}

\begin{corollary}\label{cor:b2:nvs:closedsubspace}
任一赋范空间中的有限维子空间都是闭的。
\end{corollary}

\begin{proof}
它按诱导范数是完备的
（由\cref{thm:b2:nvs:finitedim}），而完备子集都是闭的
（见\cref{def:b2:metric:complete}）。
\end{proof}

\begin{example}[借助对称性计算最佳逼近]\label{ex:b2:nvs:bestaffine}
在 $\bigl(C(\intcc{-1}{1}), \norm\cdot_\infty\bigr)$ 中，
$f(x) = \abs x$ 到（闭的、二维的）仿射函数子空间的距离
有多大？即形如 $a + bx$ 的仿射函数。由对称性，将 $a + bx$ 换成
$a - bx$ 时 $\norm{f - (a \pm bx)}_\infty$ 保持不变，且
中点 $a$ 同样不差（借助平均值的
三角不等式），故只需考虑常值函数。对常数
$a$：
\[
\norm{\abs x - a}_\infty = \max\,(1 - a,\ a)
\geq \frac12 ,
\]
在 $a = \frac12$ 处取最小，距离为 $\frac12$，
由常数 $\frac12$ 取到。注意误差曲线 $\abs x -
\frac12$：它依次在 $\pm\frac12$ 处交替取到，位于 $x = -1, 0,
1$ --- 三个符号交替的极值，对于
二参数族的最佳逼近而言。这一
\emph{等振荡}模式并非偶然，它是
最优性的标志，本章周末问题
将其推广为切比雪夫定理。
\end{example}

\begin{example}[闭子空间与稠密子空间]\label{ex:b2:nvs:closeddense}
在 $E = \bigl(C(\intcc{0}{1}), \norm\cdot_\infty\bigr)$ 中，每个
$\R_n[X]$（次数 $\leq n$ 的多项式，限制在
$\intcc01$ 上）是有限维的，因而是\emph{闭}的子空间
--- 次数 $\leq n$ 的多项式的一致极限仍属于该空间。但
所有这些子空间的并 $\R[X]$ 在 $E$ 中是\emph{稠密}的（
即魏尔斯特拉斯逼近定理，将在
\cref{ch:b2:funcseq}中证明），而稠密的真子空间
则非闭到了极致。寓意是：子空间的闭性是
有限维情形独有的一种“特权”；把一层层闭的楼层叠起来，
便可以建成一座稠密的摩天大楼。
\end{example}

\begin{theorem}[里斯]\label{thm:b2:nvs:riesz}
赋范空间 $E$ 的闭单位球是紧的，\emph{当且仅当
}$\dim E < \infty$。\index{里斯定理}
\end{theorem}

\begin{proof}
有限维时：闭且有界即可
（由\cref{thm:b2:nvs:finitedim}）。

反之，设 $\dim E = \infty$。\emph{里斯引理：}对
每个真闭子空间 $F \subsetneq E$ 和 $\varepsilon \in
\intoo{0}{1}$，存在单位向量 $x$，使 $d(x, F) \geq 1 -
\varepsilon$。证明：取 $y \notin F$，令 $\delta = d(y, F) > 0$
（因 $F$ 闭），取 $f \in F$，使 $\norm{y - f} \leq
\frac{\delta}{1 - \varepsilon}$，并令 $x = \frac{y - f}{\norm{y -
f}}$：对任意 $g \in F$，
\[
\norm{x - g} = \frac{\norm{y - (f + \norm{y-f}\,g)}}{\norm{y - f}}
\geq \frac{\delta}{\norm{y-f}} \geq 1 - \varepsilon ,
\]
分子是 $y$ 到 $F$ 中某点的距离。

现归纳地构造单位向量 $x_1, x_2, \dots$：$F_k =
\operatorname{Vect}(x_1, \dots, x_k)$ 为有限维，因而
是闭的（\cref{cor:b2:nvs:closedsubspace}）且为真子空间；由里斯引理
（取 $\varepsilon = \frac12$）得单位向量 $x_{k+1}$，使 $d(x_{k+1},
F_k) \geq \frac12$。序列满足 $\norm{x_p - x_q} \geq
\frac12$（$p \neq q$）：不存在收敛子列 --- 单位
球不紧。
\end{proof}

\begin{example}[里斯作为维数探测器]\label{ex:b2:nvs:rieszdetector}
$C(\intcc01)$ 是否为有限维？里斯无需
写出任何显式的无穷自由族即可作答：序列
$f_n(x) = x^n$ 位于闭单位球中，且对
$m > n$，$\norm{f_n - f_m}_\infty \geq f_n(x_0) - f_m(x_0) > 0$
在适当的点处成立 --- 这一点在本章周末问题
（第 16 题）中得到清晰的量化，其中某子列保持
两两距离 $\geq \frac14$。无收敛子列，故
单位球不紧，故 $\dim C(\intcc01) = \infty$，
由\cref{thm:b2:nvs:riesz}可知。单位球的紧性是一种
完美的二分：在有限维中成立，在无穷维中
失效，两者之间没有中间地带 --- 仅凭几何即可
读出空间的维数类型。
\end{example}

\section{巴拿赫空间}

\begin{definition}\label{def:b2:nvs:banach}
称\emph{巴拿赫空间}\index{巴拿赫空间}为完备
赋范空间。例：每个有限维赋范空间
（\cref{thm:b2:nvs:finitedim}）；$\bigl(C(\intcc{a}{b}),
\norm\cdot_\infty\bigr)$（\cref{thm:b2:metric:rncomplete}）；
$\mathcal{L}_c(E, F)$，$F$ 为巴拿赫空间（证明方式与
连续函数的情形相同）。反例：$\bigl(C(\intcc{0}{1}),
\norm\cdot_1\bigr)$（\cref{exo:b2:nvs:7}）。
\end{definition}

\begin{example}[积分的算子范数]\label{ex:b2:nvs:integrationnorm}
在 $\bigl(C(\intcc01), \norm\cdot_\infty\bigr)$ 上，令 $T(f)(x) =
\int_0^x f(t)\,\dd t$（这是一个自同态：$T(f)$ 连续）。
运用 \cref{met:b2:nvs:opnorm}。上界：
\[
\abs{T(f)(x)} \leq \int_0^x\abs f \leq x\,\norm f_\infty \leq
\norm f_\infty ,
\]
故 $\vertiii T \leq 1$。取 $f \equiv 1$ 为见证，给出 $T(f)(x) =
x$，且 $\norm{T(f)}_\infty = 1 = \norm f_\infty$：取到，
$\vertiii T = 1$。但注意 $\vertiii{T^2} = \frac12 \neq
\vertiii T^2$：事实上 $T^2(f)(x) = \int_0^x(x - t)f(t)\dd t$ 有
$\abs{T^2(f)(x)} \leq \frac{x^2}2\norm f_\infty$，同样在 $f \equiv 1$ 时取到，
一般地 $\vertiii{T^n} = \frac1{n!}$ ---
次乘性界 $\vertiii T^n = 1$ 相差一个阶乘。这正是
迭代技巧所揭示的现象，它在
\cref{ch:b2:metric}的周末问题中被转化为线性微分方程
的全局可解性。
\end{example}

\begin{theorem}[巴拿赫空间中的绝对收敛]\label{thm:b2:nvs:absoluteconvergence}
在巴拿赫空间中，若 $\sum \norm{u_n} < \infty$，则 $\sum u_n$
收敛，且 $\norm{\sum u_n} \leq \sum\norm{u_n}$。（
赋范空间中级数的完整理论见\cref{ch:b2:series}。）
\end{theorem}

\begin{proof}
部分和 $S_N$：对 $q > p$，$\norm{S_q - S_p} \leq
\sum_{n=p+1}^{q}\norm{u_n}$，它趋于 $0$（柯西准则，
用于由范数构成的实级数）：$(S_N)$ 为柯西列，因而收敛。
由有限和的三角不等式取极限，即得
该不等式。
\end{proof}

\begin{example}[矩阵指数初探]\label{ex:b2:nvs:matrixexp}
\emph{矩阵指数}\index{矩阵指数}：
$\mathcal{M}_n(K)$ 配以任一次乘性范数（$\vertiii{AB}
\leq \vertiii A \vertiii B$）时是巴拿赫空间（有限维）。于是
对每个 $A$，
\[
\eu^A = \sum_{k=0}^{\infty} \frac{A^k}{k!}
\]
绝对收敛（$\vertiii{A^k/k!} \leq \vertiii A^k /k!$，
可和）：良定义。\cref{ch:b2:diffeq}中系统地
利用了这一点。
\end{example}

\begin{example}[终止的诺伊曼级数]\label{ex:b2:nvs:neumannfinite}
取 $A = \begin{psmallmatrix}0 & \frac12\\ 0 &
0\end{psmallmatrix}$：$\vertiii A < 1$，对由
\cref{ex:b2:nvs:examples}中范数构造的任一算子范数均成立，且 $A^2 = 0$，故
几何级数于是退化为有限和：
\[
(I - A)^{-1} = \sum_{k \geq 0} A^k = I + A =
\begin{pmatrix}1 & \tfrac12\\ 0 & 1\end{pmatrix},
\]
由 $(I - A)(I + A) = I - A^2 = I$ 验证。幂零性恰如
截断此级数，恰如它在
\cref{ch:b2:reduction}中截断指数级数那样；这个例子还校准了
期望：诺伊曼逆在
一般情形下是无穷级数，仅当扰动为
幂零时才是多项式，且 $N$ 项之后的误差总被
几何尾项 $\vertiii A^{N+1}/(1 - \vertiii A)$ 所控制。
\end{example}

\begin{example}[旋转生成元的指数]\label{ex:b2:nvs:rotationexp}
设 $A = \begin{psmallmatrix}0 & -\theta\\ \theta &
0\end{psmallmatrix}$。则 $A^2 = -\theta^2 I$，故各次幂
以四为周期循环，级数拆为偶次项与奇次项
两部分：
\[
\eu^{A} = \sum_{k}\frac{A^k}{k!}
= \Bigl(\sum_{j}\frac{(-1)^j\theta^{2j}}{(2j)!}\Bigr) I
+ \Bigl(\sum_{j}\frac{(-1)^j\theta^{2j+1}}{(2j+1)!}\Bigr)
\frac{A}{\theta}
= \begin{pmatrix}
\cos\theta & -\sin\theta\\
\sin\theta & \cos\theta
\end{pmatrix},
\]
各项的重排均由绝对收敛许可。反对称生成元的
指数是一个旋转 ---
此处仅由级数算出，早于三章之后的
微分方程 $x' = Ax$（\cref{ch:b2:diffeq}）解释了
\emph{为何}：$\eu^{tA}$ 是匀速圆周运动。本例的
结论是：矩阵级数之间的恒等式证明方式与
数量级数完全相同，只要一个次乘范数能保证绝对
收敛。
\end{example}

\begin{example}[上确界范数即一致收敛：对照表]\label{ex:b2:nvs:uniformdictionary}
该断言 $\norm{f_n - f}_\infty \to 0$ \emph{正是}一致
收敛：一个数 $\sup_x\abs{f_n(x) - f(x)}$ 同时界定
了每一点处的误差。对照表在
$f_n(x) = x^n$ 于 $\intcc{0}{1}$ 上的应用：逐点而言，$f_n
\to 0$ 于 $\intco{0}{1}$ 上，且 $f_n(1) = 1$；就范数而言，
$\norm{f_n - 0}_\infty = 1 \not\to 0$，事实上逐点
极限不连续，因而它不可能是
$\norm\cdot_\infty$ 意义下 $C(\intcc01)$ 中的极限（该空间对一致极限封闭，
见\cref{thm:b2:metric:rncomplete}）。在
$\intcc{0}{a}$ 上，$a < 1$：$\norm{f_n}_\infty = a^n \to 0$ ---
经缩小定义域即可恢复一致收敛。凡
在\cref{ch:b2:funcseq}中出现的收敛陈述，皆
是关于这一范数的陈述；牢记此对照表，可使
该章的难度减半。
\end{example}

\begin{remark}[常见误区]\label{rem:b2:nvs:pitfalls}
(i) 算子范数同时依赖于所取的\emph{两个}范数：同一
矩阵的 $\vertiii\cdot_\infty$ 由行和给出
（\cref{exo:b2:nvs:4}），另一个 $\vertiii\cdot_1$ 由列
和给出；笼统地谈论矩阵的“那个”范数，
不指明底层范数便毫无意义。(ii) $\vertiii{AB} \leq
\vertiii A\,\vertiii B$ 是一个不等式，通常为严格的 ---
幂的衰减可以远快于界 $\vertiii A^k$
所暗示的，这正是适配范数的全部意义（本章
周末问题第 22 题）。(iii) “线性蕴含
连续”仅是有限维空间的特权：求导
在多项式上既线性又无界
（\cref{ex:b2:nvs:operatornorms}）。(iv) $\sum u_n$ 的绝对收敛
只有在空间完备时才有助益
（\cref{exo:b2:series:9} 构造了反例）。(v) 在
无限维中，单位球上的上确界是真正的
上确界，不可想当然地认为它一定可取到
（\cref{exo:b2:nvs:8}）。
\end{remark}

\begin{remark}[本册内的几个视角]\label{rem:b2:nvs:perspectives}
三处关联现已确定。与\cref{ch:b2:series}相关：在
巴拿赫空间中，绝对收敛的级数必收敛，因而
算子的几何级数与指数级数成为日常
工具：求 $I - A$ 的逆，定义 $\eu^{A}$
（\cref{ex:b2:series:neumann}）。与
\cref{ch:b2:funcseq}和\cref{ch:b2:powerseries}相关：函数列与幂级数的收敛
性，即是在
$\bigl(C, \norm\cdot_\infty\bigr)$ 中的收敛
（\cref{ex:b2:nvs:uniformdictionary}），而收敛半径
是关于哪些几何级数占优的陈述。与
\cref{ch:b2:fourier}相关：范数
$\norm\cdot_2$ 与 $\norm\cdot_\infty$ 在
$C(\intcc{0}{1})$ 上确实不一致（\cref{ex:b2:nvs:sqrtnxn}），这恰好解释了
为何傅里叶级数的均方收敛与一致
收敛是两条不同的定理，各付出不同的
代价。
\end{remark}

\begin{remark}[本章的应用]
算子范数与几何级数支撑着
在\cref{ch:b2:diffcalc}中的扰动论证（逆函数定理）以及
在\cref{ch:b2:diffeq}中的矩阵指数；范数等价性
默默授权了每一个“任取你喜爱的范数”的论证，
在\cref{ch:b2:funcseq}及以后的各章中；而
里斯定理中有限与无限的分野 --- 经本章周末问题定量化 ---
正是本科三年级分册需要新工具的原因
（弱收敛、阿尔泽拉–阿斯科利定理、希尔伯特空间投影），
而本册仍可在此提取
收敛子列。
\end{remark}

\section{练习}

\begin{exercise}[$\star$]\label{exo:b2:nvs:1}
在 $\R^2$ 上画出 $\norm\cdot_1$、$\norm\cdot_2$、
$\norm\cdot_\infty$ 的单位球，并证明不等式 $\norm x_\infty
\leq \norm x_2 \leq \norm x_1 \leq 2\norm x_\infty$，并
求出维数 $2$ 时的最佳常数。
\end{exercise}

\begin{exercise}[$\star$]\label{exo:b2:nvs:2}
判断 $N(f) = \abs{f(0)} + \norm{f'}_\infty$ 是否为
$C^1(\intcc{0}{1})$ 上的范数？将它与 $\norm{f}_\infty$ 比较，其中一个
不等式成立，另一个不成立（请给出反例）。
\end{exercise}

\begin{exercise}[$\star$]\label{exo:b2:nvs:3}
计算 $u(f) = \int_0^1 f(t)\,\eu^t\,\dd t$ 的算子范数，
在 $\bigl(C(\intcc{0}{1}), \norm\cdot_\infty\bigr) \to \R$ 上，以及
平移算子 $S(x_1, x_2, \dots, x_n) = (x_2, \dots, x_n, 0)$ 在
$(K^n, \norm\cdot_\infty)$ 上的算子范数。
\end{exercise}

\begin{exercise}[$\star\star$]\label{exo:b2:nvs:4}
在 $(\R^n, \norm\cdot_\infty)$ 上，证明矩阵 $A$ 的算子范数
为 $\vertiii A_\infty = \max_i \sum_j \abs{a_{ij}}$（即
最大绝对行和）。并求矩阵 $\begin{pmatrix} 1 & -2\\
3 & 1\end{pmatrix}$ 的算子范数。
\end{exercise}

\begin{exercise}[$\star\star$]\label{exo:b2:nvs:5}
证明 $GL_n(K)$ 是 $\mathcal{M}_n(K)$ 中的开集，并证明 $A
\mapsto A^{-1}$ 在其上连续。\emph{提示：证明开性时，若
$\vertiii H < \frac{1}{\vertiii{A^{-1}}}$ 则 $A + H = A(I +
A^{-1}H)$，且 $\vertiii{A^{-1}H} < 1$，则 $I + B$ 在
$\vertiii B < 1$ 时可逆，由几何级数
（\cref{thm:b2:nvs:absoluteconvergence}）；对连续性，
估计 $(A+H)^{-1} - A^{-1}$，利用同一级数。}
\end{exercise}

\begin{exercise}[$\star\star$]\label{exo:b2:nvs:6}
设 $\varphi$ 为赋范空间 $E$ 上的线性型。证明
$\varphi$ 连续当且仅当 $\ker\varphi$ 是闭集。
\emph{（若 $\ker\varphi$ 为闭集且 $\varphi \neq 0$，取 $a$
使 $\varphi(a) = 1$，并取 $r > 0$ 使 $B(a, r) \cap \ker\varphi =
\emptyset$。要推出 $\abs{\varphi(h)} \leq \frac{1}{r}\norm h$，
可对 $a - \frac{h}{\varphi(h)}$ 作伸缩论证。）}
\end{exercise}

\begin{exercise}[$\star\star$]\label{exo:b2:nvs:7}
证明 $\bigl(C(\intcc{0}{1}), \norm\cdot_1\bigr)$ 不是
完备的：证明函数列 $f_n$，即从 $0$ 到
$1$ 的仿射斜坡，定义在 $\bigl[\frac12 - \frac1n, \frac12\bigr]$ 上（左侧取值 $0$，
右侧取值 $1$），证明其为柯西列，且无连续的
$\norm\cdot_1$ 极限。
\end{exercise}

\begin{exercise}[$\star\star\star$]\label{exo:b2:nvs:8}
在 $E = C(\intcc{0}{1})$ 上赋予 $\norm\cdot_\infty$，考虑
\[
\varphi(f) = \sum_{n \geq 1} (-1)^n\, 2^{-n} f\bigl(\tfrac1n\bigr).
\]
证明 $\varphi$ 是良定义的连续线性型，且
$\vertiii\varphi = 1$，但定义 $\vertiii\varphi$ 的上确界
在闭单位球上\emph{不能达到}。
\emph{（上界由三角不等式给出。范数 $= 1$：构造
连续的 $f_K$，满足 $\norm{f_K}_\infty \leq 1$ 且 $f_K(\frac1n)
= (-1)^n$，当 $n \leq K$ 时 --- 点 $\frac1n$ 彼此孤立。
不能达到：若取等，则 $f(\frac1n) =
(-1)^n$ 对每个 $n$ 成立，与 $f$ 在 $0$ 处连续相矛盾
因 $\frac1n \to 0$。）}
\end{exercise}

\begin{exercise}[$\star\star\star$]\label{exo:b2:nvs:9}
设 $E$ 为赋范空间，其闭单位球为紧集。
不引用\cref{thm:b2:nvs:riesz}，重新推导
每个有界序列都有收敛子列，并证明 $E$ 上的每个
线性型连续当且仅当 $\dim E < \infty$。
\emph{（对无限维情形，构造一个不连续的线性型：
先在一个线性无关的规范化序列上任意定义它，
再将其延拓 --- 在本阶段不加证明地承认代数补
的存在性。）}
\end{exercise}

\begin{exercise}[$\star\star$]\label{exo:b2:nvs:10}
在 $C(\intcc{0}{1})$ 上，证明 $\norm f_1 \leq \norm f_2 \leq \norm
f_\infty$，\emph{（对第一个用柯西--施瓦茨不等式）}，并
通过函数族 $f_n(x) = x^n$ 表明两个不等式都不能反过来，
即在相差常数倍的意义下，三个范数两两不等价。
\end{exercise}

\begin{exercise}[$\star\star$]\label{exo:b2:nvs:11}
（到超平面的距离）设 $\varphi$ 为非零连续的
赋范空间 $E$ 上的线性型，证明
\[
d\bigl(x, \ker\varphi\bigr) =
\frac{\abs{\varphi(x)}}{\vertiii\varphi}
\qquad (x \in E),
\]
并在\cref{exo:b2:nvs:8}中验证，下确界未必
能被超平面上的任何一点达到。
\end{exercise}

\begin{exercise}[$\star\star\star$]\label{exo:b2:nvs:12}
在 $E = \R[X]$（全体多项式）上，令 $N_1(P) =
\sup_{\intcc{0}{1}}\abs P$ 以及 $N_2(P) =
\sup_{\intcc{0}{2}}\abs P$。证明 $N_1 \leq N_2$，但
$N_1$ 与 $N_2$ \emph{不}等价；由此推出
恒等映射 $(E, N_2) \to (E, N_1)$ 是连续线性双射，
其逆映射不连续。最后证明 $(E, N_1)$
不完备 \emph{（$\eu^x$ 的泰勒部分和）}。以上三种
现象在有限维中都不可能出现 --- 请说明理由。
\end{exercise}

\section{问题：最佳逼近与切比雪夫定理}

给定次数的多项式能把一个函数逼近到何种程度？
哪个多项式逼近得最好？就存在性而言，
这一问题属于本章：有限维空间中的紧性
保证了最佳逼近的存在。就具体构造而言，
有一个非平凡的情形可以仅用初等方法完全求解 ---
在所有次数为 $n$ 的\emph{首一}多项式中，
在 $\intcc{-1}{1}$ 上的上确界范数最小者，即（经归一化的）切比雪夫多项式，
其范数为 $2^{1-n}$：\emph{切比雪夫极值
定理}。本题两方面都予以证明，然后度量
紧性在无限维中失效的程度：
$C(\intcc{0}{1})$ 的单位球包含无穷多个
两两距离为 $1$ 的点构成的“星座”。

\begin{problem}[{周末问题 --- 切比雪夫极值定理
与单位球的几何}]
\label{pb:b2:nvs:1}
不带下标的范数均指所示区间上的上确界范数。

\medskip
\noindent\textbf{第一部分 --- 赋范空间中的最佳
逼近。}
\begin{enumerate}
  \item 设 $F$ 为赋范空间
        $E$ 的有限维子空间，$x \in E$。证明距离 $d(x,
        F) = \inf_{f \in F}\norm{x - f}$ 的下确界\emph{可以取到}
        \emph{（化为 $F$ 中一个有界闭子集上的问题，并利用
        \cref{thm:b2:nvs:finitedim}）}。
  \item 称范数为\emph{严格凸}的，若 $\norm u = \norm v
        = 1$ 且 $u \neq v$，则 $\bigl\Vert\frac{u +
        v}2\bigr\Vert < 1$。证明 $\norm\cdot_2$ 在 $\R^n$
        上是严格凸的\emph{（利用平行四边形恒等式）}，并
        说明 $\norm\cdot_1$ 与 $\norm\cdot_\infty$ 在
        $n \geq 2$ 时都不是严格凸的。
  \item 证明：对严格凸范数，第 1 小题中的最佳
        逼近是\emph{唯一}的。
  \item 在 $(\R^2, \norm\cdot_\infty)$ 中，求 $x = (0, 1)$ 关于
        直线 $F =
        \operatorname{Vect}\bigl((1,0)\bigr)$ 的全部最佳逼近，它们
        构成一个区间。
  \item 在 $\bigl(C(\intcc{a}{b}), \norm\cdot_\infty\bigr)$ 中，
        证明 $f$ 关于
        \emph{常值函数}的最佳逼近唯一，等于 $c^* = \frac{\max f
        + \min f}{2}$，距离为 $\frac{\max f - \min
        f}{2}$；对 $f(x) = x^2$ 在
        $\intcc{0}{1}$ 上，求出这两个量。
\end{enumerate}

\medskip
\noindent\textbf{第二部分 --- 切比雪夫多项式。}
\begin{enumerate}[resume]
  \item 证明恰存在一个多项式 $T_n$，满足
        $T_n(\cos\theta) = \cos n\theta$ 对所有 $\theta$ 成立，
        \emph{（递推式 $T_{n+1} = 2XT_n - T_{n-1}$ 由
        余弦加法公式推出）}，$\deg T_n = n$，并且
        其首项系数为 $2^{n-1}$，其中 $n \geq 1$。
  \item 证明 $\abs{T_n} \leq 1$ 在 $\intcc{-1}{1}$ 上成立，且
        $T_n(\eta_k) = (-1)^k$，在 $n + 1$ 个点 $\eta_k =
        \cos\frac{k\pi}{n}$（$k = 0, \dots, n$）处成立，且
        $T_n$ 的零点恰为 $n$ 个点
        $\cos\frac{(2k-1)\pi}{2n}$，它们与 $\eta_k$ 相互交错。
  \item 计算 $T_2, T_3, T_4$，并通过直接求值验证
        $T_3$ 的交错性（取 $\eta_0, \dots, \eta_3 = 1, \frac12, -\frac12,
        -1$ 处的值）。
  \item 对 $\abs x \geq 1$，证明
        \[
        T_n(x) = \frac{\bigl(x + \sqrt{x^2 - 1}\bigr)^n +
        \bigl(x - \sqrt{x^2 - 1}\bigr)^n}{2},
        \]
        并由此推出 $T_n(x) \sim \frac12\bigl(x + \sqrt{x^2 -
        1}\bigr)^n \to \infty$，对固定的 $x > 1$ 呈几何级数增长。
  \item 证明复合律 $T_m \circ T_n = T_{mn}$
        \emph{（在 $\intcc{-1}{1}$ 上验证，并利用
        多项式的刚性）}。
\end{enumerate}

\medskip
\noindent\textbf{第三部分 --- 切比雪夫极值定理。}
记 $Q_n = 2^{1-n}T_n$，由第 6 小题可知它是首一的。
\begin{enumerate}[resume]
  \item 设 $P$ 为次数 $n \geq 1$ 的首一多项式，且
        $\sup_{\intcc{-1}{1}}\abs P < 2^{1-n}$。令 $D
        = Q_n - P$，在点 $\eta_k$ 处求值并计数
        符号变化，导出矛盾。于是得出：
        \[
        \sup_{\intcc{-1}{1}}\abs P \;\geq\; 2^{1-n}
        \qquad\text{对每个首一的 } P \text{，次数为 } n.
        \]
  \item （等号情形）设 $\sup_{\intcc{-1}{1}}\abs P =
        2^{1-n}$，其中 $P$ 为次数 $n$ 的首一多项式，并令 $D = Q_n
        - P \neq 0$。证明 $(-1)^kD(\eta_k) \geq 0$ 对所有 $k$ 成立；
        证明这 $n$ 个区间中的每一个
        $\intcc{\eta_{k}}{\eta_{k-1}}$ 都含有 $D$ 的一个零点，
        并且两个相邻区间共享的零点是一个
        \emph{内点} $\eta_k$，且此处 $D' = 0$ 也成立。
        于是 $D$ 按重数计恰有 $n$ 个零点，
        因而 $D = 0$，即极小元恰为
        $Q_n$ --- \emph{切比雪夫极值定理}。
  \item 将定理改述为距离形式：在 $\intcc{-1}{1}$ 上，
        \[
        d_\infty\bigl(X^n,\ \R_{n-1}[X]\bigr) = 2^{1-n},
        \]
        其唯一的最佳逼近为 $X^n - Q_n$，并
        通过仿射代换 $x = \frac{1+t}2$ 证明在
        $\intcc{0}{1}$ 上距离变为 $2^{1-2n}$。
  \item （最优插值节点）对 $n$ 个节点 $x_1, \dots,
        x_n \in \intcc{-1}{1}$，节点多项式
        $\omega(x) = \prod_i(x - x_i)$ 是次数为 $n$ 的首一多项式。
        由第 12 小题推出，何种节点选取使
        $\sup_{\intcc{-1}{1}}\abs\omega$ 最小（该量是经典插值误差界
        中依赖于节点的因子），并
        求出其最小值。
  \item 用手算验证定理在 $n = 2$ 的情形（直接求
        $\inf_c \sup_{\intcc{-1}{1}}\abs{x^2 - c}$），
        并数值计算第 13 小题中的距离，
        即在 $\intcc{0}{1}$ 上取 $n = 10$。该距离的大小说明
        $x^{10}$ 的图像有怎样的特点？
\end{enumerate}

\medskip
\noindent\textbf{第四部分 --- $C(\intcc{0}{1})$
的单位球。}
\begin{enumerate}[resume]
  \item 设 $g_k(x) = x^{2^k}$。证明 $\norm{g_k}_\infty = 1$，且
        $\norm{g_k - g_j}_\infty \geq \frac14$，对 $j > k$ 成立
        \emph{（在使 $x^{2^k} =
        \frac12$ 的点处求值）}：得到一个显式有界序列，它没有
        收敛子列，即闭单位球不是
        紧的，可直接验证。
  \item （里斯引理的加强）设 $F$ 为
        \emph{有限维}的真子空间，其所在赋范空间为
        $E$。利用第 1 小题，取一个单位向量 $x$，
        使 $d(x, F) = 1$ 恰好成立，而不仅是 $\geq 1 -
        \varepsilon$，如\cref{thm:b2:nvs:riesz}的引理中所述。
  \item 由此推出：每个无穷维赋范空间中都有一列
        单位向量，两两距离均 $\geq 1$，
        并由此重新导出里斯定理。
  \item 在 $C(\intcc{0}{1})$ 中显式给出这样一个“星座”：
        取帐篷函数 $h_n$，其支集为
        $\bigl[\frac1{n+1}, \frac1n\bigr]$，峰值为 $1$。
        验证 $\norm{h_n} = 1$，$\norm{h_n - h_m} = 1$（当 $n
        \neq m$ 时），并注意 $h_n \to 0$ 逐点收敛，但并非
        一致收敛。
  \item （全有界性不成立）证明
        $C(\intcc{0}{1})$ 的闭单位球不能被有限个
        半径为 $\frac13$ 的球覆盖，\emph{（每个这样的球至多含有一个
        $h_n$）}——与
        \cref{thm:b2:metric:borellebesgue}证明中的全有界性步骤
        形成对照。
\end{enumerate}

\medskip
\noindent\textbf{第五部分 --- 矩阵上的范数应用，以及
综合}
\begin{enumerate}[resume]
  \item 证明：若 $\lambda$ 是 $A \in
        \mathcal{M}_n(\C)$ 的特征值，则 $\abs\lambda \leq
        \vertiii A$ 对每个算子范数成立；应用
        \cref{exo:b2:nvs:4}，估计矩阵
        $\begin{psmallmatrix}1 & -2\\ 3 & 1\end{psmallmatrix}$
        的特征值的界，并与它们的真实模长比较。
  \item （适配范数）设 $A$ 可对角化，$A =
        P\,\mathrm{diag}(\lambda_1, \dots, \lambda_n)\,P^{-1}$。
        证明 $N_P(x) = \norm{P^{-1}x}_\infty$ 是一个范数，其
        算子范数满足 $\vertiii A_{N_P} =
        \max_i\abs{\lambda_i}$。
  \item 由此推出：对可对角化的 $A$，$A^k \to 0$ 当且仅当
        所有特征值满足 $\abs{\lambda_i} < 1$ ---
        范数的等价性使结论
        与所用范数无关。以 $A =
        \frac14\begin{psmallmatrix}1 & 2\\ 2 &
        1\end{psmallmatrix}$ 为例验证。
  \item （等价常数无限增大）在 $\R_n[X]$ 上比较
        $N_c(P) = \max_k \abs{a_k}$（按系数）与
        $\norm{P}_{\intcc{0}{1}}$：二者都是范数，故
        对每个固定的 $n$ 它们等价；但利用
        第 13 题中在 $\intcc{0}{1}$ 上的首一极小化多项式，
        证明：记 $C_n$ 为使 $N_c \leq
        C_n\norm\cdot_{\intcc{0}{1}}$ 成立的最佳常数，则 $C_n \geq
        2^{2n-1}$。用一句话说明为何“所有范数都等价”
        在无限维中失效。
  \item （综合）各用一句话说明：有限维球的紧性在何处起作用
        （第 1、12 题）；严格凸性
        支配着什么；第 18--19 题这一组题破坏了什么；
        以及第 24 题如何定量刻画这种失效。
        点明其顶峰（切比雪夫
        极值定理），并说明最佳逼近的现代归宿
        （希尔伯特空间上的投影定理，
        见本科三年级分册，其中完备性取代了
        紧性）。
\end{enumerate}
\end{problem}
